\documentclass[12pt,a4paper]{article}
\usepackage[T1]{fontenc}
\usepackage{amsmath,amssymb,mathrsfs,tikz,times,pifont}
\usepackage{enumitem}
\usepackage{multicol}
\usepackage{lmodern}
\usetikzlibrary{trees}
\newcommand\circitem[1]{%
\tikz[baseline=(char.base)]{
\node[circle,draw=gray, fill=red!55,
minimum size=1.2em,inner sep=0] (char) {#1};}}
\newcommand\boxitem[1]{%
\tikz[baseline=(char.base)]{
\node[fill=cyan,
minimum size=1.2em,inner sep=0] (char) {#1};}}
\setlist[enumerate,1]{label=\protect\circitem{\arabic*}}
\setlist[enumerate,2]{label=\protect\boxitem{\alph*}}
%%%::::::by chnini ameur :::::::%%%
\everymath{\displaystyle}
\usepackage[left=1cm,right=1cm,top=1cm,bottom=1.7cm]{geometry}
\usepackage[colorlinks=true, linkcolor=blue, urlcolor=blue, citecolor=blue]{hyperref}
\usepackage{array,multirow}
\usepackage[most]{tcolorbox}
\usepackage{varwidth}
\usepackage{float} %pour utiliser l'option [H] qui force l'image à apparaître exactement à l'endroit où elle est placée dans le code.
\tcbuselibrary{skins,hooks}
\usetikzlibrary{patterns}
%%%::::::by chnini ameur :::::::%%%
\newtcolorbox{exa}[2][]{enhanced,breakable,before skip=2mm,after skip=5mm,
colback=yellow!20!white,colframe=black!20!blue,boxrule=0.5mm,
attach boxed title to top left ={xshift=0.6cm,yshift*=1mm-\tcboxedtitleheight},
fonttitle=\bfseries,
title={#2},#1,
% varwidth boxed title*=-3cm,
boxed title style={frame code={
\path[fill=tcbcolback!30!black]
([yshift=-1mm,xshift=-1mm]frame.north west)
arc[start angle=0,end angle=180,radius=1mm]
([yshift=-1mm,xshift=1mm]frame.north east)
arc[start angle=180,end angle=0,radius=1mm];
\path[left color=tcbcolback!60!black,right color = tcbcolback!60!black,
middle color = tcbcolback!80!black]
([xshift=-2mm]frame.north west) -- ([xshift=2mm]frame.north east)
[rounded corners=1mm]-- ([xshift=1mm,yshift=-1mm]frame.north east)
-- (frame.south east) -- (frame.south west)
-- ([xshift=-1mm,yshift=-1mm]frame.north west)
[sharp corners]-- cycle;
},interior engine=empty,
},interior style={top color=yellow!5}}
%%%%%%%%%%%%%%%%%%%%%%%
\usepackage{fancyhdr}
\usepackage{eso-pic}         % Pour ajouter des éléments en arrière-plan
% Commande pour ajouter du texte en arrière-plan
\usepackage{tkz-tab}
\AddToShipoutPicture{
    \AtTextCenter{%
        \makebox[0pt]{\rotatebox{80}{\textcolor[gray]{0.7}{\fontsize{5cm}{5cm}\selectfont PGB}}}
    }
}
\usepackage{lastpage}
\fancyhf{}
\pagestyle{fancy}
\renewcommand{\footrulewidth}{1pt}
\renewcommand{\headrulewidth}{0pt}
\renewcommand{\footruleskip}{10pt}
\fancyfoot[R]{
\color{blue}\ding{45}\ \textbf{2025}
}
\fancyfoot[L]{
\color{blue}\ding{45}\ \textbf{Prof:M. BA}
}
\cfoot{\bf
\thepage /
\pageref{LastPage}}
% Création du compteur pour les exercices
\newcounter{exercice}
\renewcommand{\theexercice}{\arabic{exercice}}  % Définit l'affichage du compteur en chiffres arabes

% Définir la commande \exo
\newcommand{\exo}{\refstepcounter{exercice}\textbf{Exercice \theexercice} }
\begin{document}
\renewcommand{\arraystretch}{1.5}
\renewcommand{\arrayrulewidth}{1.2pt}
\begin{tikzpicture}[overlay,remember picture]
    \node[draw=blue,line width=1.2pt,fill=purple,text=blue,inner sep=3mm,rounded corners,pattern=dots]at ([yshift=-2.5cm]current page.north) {\begingroup\setlength{\fboxsep}{0pt}\colorbox{white}{\begin{tabular}{|*1{>{\centering \arraybackslash}p{0.28\textwidth}} |*2{>{\centering \arraybackslash}p{0.2\textwidth}|} *1{>{\centering \arraybackslash}p{0.19\textwidth}|} }
                \hline
                \multicolumn{3}{|c|}{$\diamond$$\diamond$$\diamond$\ \textbf{Lycée de Dindéfélo}\ $\diamond$$\diamond$$\diamond$ } & \textbf{A.S. : 2026/2027}                                              \\ \hline
                \textbf{Matière: Mathématiques}                                                                                    & \textbf{Niveau : }\textbf{TL} & \textbf{Date: 05/11/2026} & \textbf{} \\ \hline
                \multicolumn{4}{|c|}{\parbox[c]{10cm}{\begin{center}
                                                                  \textbf{{\Large\sffamily Td Polynomes}}
                                                              \end{center}}}                                                                                                        \\ \hline
            \end{tabular}}\endgroup};
\end{tikzpicture}

\vspace{3cm}

\fbox{\textbf{\exo}}
Le réel $\alpha$ est une racine d'un polynôme
\begin{enumerate}
    \item $\alpha$ est nul
    \item $P(\alpha) \neq 0$
    \item $P(x)$ est divisible par $(x - \alpha)$
\end{enumerate}

\fbox{\textbf{\exo}(5 points)}

On pose $P(x) = 2x^3 + 11x^2 + 4x - 5$

\begin{enumerate}
    \item Montrer que $-5$ est une racine de $P(x)$. \hfill (1pt)
    \item Factoriser $P(x)$. \hfill (1pt)
    \item Résoudre dans $\mathbb{R}$ l'équation $P(x) = 0$. \hfill (1pt)
    \item Résoudre dans $\mathbb{R}$ l'inéquation $P(x) \leq 0$. \hfill (1pt)
    \item Déduire de la question 3) l'ensemble solution de l'équation $2(x - 2)^3 + 11(x - 2)^2 + 4(x - 2) - 5 = 0$.
\end{enumerate}

\fbox{\textbf{\exo}(10 points)}

\textbf{A.} Soit $P$ le polynôme défini par : $P(x) = 2x^3 + \mathbf{a}x^2 + 12x + \mathbf{b}$ où $a$ et $b$ sont des réels
\begin{enumerate}
    \item Trouver les réels $\mathbf{a}$ et $\mathbf{b}$ tels que $P(2) = 5$ et $P(-1) = -16$ \hfill \textbf{(2pts)}
\end{enumerate}

\textbf{B.} On considère le polynôme $P$ défini par : $P(x) = 2x^3 - 11x^2 + 12x + 9$.
\begin{enumerate}
    \item Montrer qu'il existe un polynôme $Q$, dont on précisera le degré, tel que : \\
    $\forall x \in \mathbb{IR}, P(x) = (x - 3)Q(x)$. \hfill \textbf{(1pt)}
    \item Déterminer $Q(x)$. \hfill \textbf{(1pt)}
    \item Résoudre dans $\mathbb{IR}$ : a) $P(x) = 0$ ; b) $P(1 - 4x) = 0$ ; c) $P(x) \leq 0$ \hfill \textbf{(1.5pts)}
\end{enumerate}

\textbf{C.} Soit $F$ la fraction rationnelle définie par : $F(x) = \frac{P(x)}{x^2 - 9}$.
\begin{enumerate}
    \item Déterminer le domaine de définition $D_F$ de $F$. \hfill \textbf{(1pt)}
    \item Simplifier $F(x)$ dans $D_F$. \hfill \textbf{(1pt)}
    \item Résoudre dans $\mathbb{IR}$ : a) $F(x) = 0$ ; b) $F(x) \geq 0$. \hfill \textbf{(1pt)}
    \item Déterminer les nombres réels $a$, $b$ et $c$ tels que : $\forall x \in D_F, F(x) = ax + b + \frac{c}{x + 3}$. \hfill \textbf{(1.5pts)}
\end{enumerate}

\fbox{\textbf{\exo}(12 points)}

Considérons le polynôme $P(x) = -3x^3 + ax^2 - 9x + b$ où $a$ et $b$ sont des nombres réels .
\begin{enumerate}
\item Déterminer $a$ et $b$ sachant que $P\left(\frac{1}{3}\right) = 0$ et $P(-1) = 24$ . \hfill \textbf{(2pts)}
\item On pose $a = 10$ et $b = 2$ .
\begin{enumerate}
    \item Calculer $P(2)$ puis conclure . \hfill \textbf{(0,5pt)}
    \item Factoriser de façon complète le polynôme $P(x)$ . \hfill \textbf{(2pts)}
    \item Résoudre dans $\mathbb{IR}$, l'équation $P(x) = 0$ . En déduire la solution de $P\left(\frac{x+1}{2x}\right) = 0$ . \hfill \textbf{(1,5pt)}
    \item Résoudre dans $\mathbb{IR}$, l'inéquation $P(x) > 0$ . \hfill \textbf{(1pt)}
\end{enumerate}
\item Soit la fonction $f$ définie par : $f(x) = \frac{P(x)}{3x^2 + 2x - 1}$ .

\begin{enumerate}
    \item Déterminer le domaine de définition $D_f$ de $f$ . \hfill \textbf{(2pts)}
    \item Simplifier $f(x)$ sur $D_f$ . \hfill \textbf{(1pt)}
    \item Etudier le signe de $f(x)$ sur $D_f$ à l'aide d'un tableau de signe . \hfill \textbf{(1,5pt)}
    \item En déduire la solution de l'inéquation $f(x) \leq 0$ . \hfill \textbf{(0,5pt)}
\end{enumerate}
\end{enumerate}
\fbox{\textbf{\exo}(07 points)}

\vspace{0.2cm}

\noindent On considère le polynôme P défini par $p(x) = 2x^3 - 7x^2 - 5x + 4$

\begin{enumerate}
    \item Montrer que 4 est une racine de P . \hfill \textbf{(0,5pt)}
    \item Factoriser $P(x)$ \hfill \textbf{(1,5 pts)}
    \item Résoudre $P(x) = 0$ et $P(x) \geq 0$ \hfill \textbf{(3pts)}
    \item En déduire les solutions de :
    \begin{enumerate}
        \item $2x^6 - 7x^4 - 5x^2 + 4 = 0$ \hfill \textbf{(1pt)}
        \item $P(-2x + 1) = 0$ \hfill \textbf{(1pt)}
    \end{enumerate}
\end{enumerate}

\fbox{\textbf{\exo}(09 points)}

\noindent On considère le polynôme $Q$ définie par: $Q(x) = (x + 3)(x - 1)(-x^2 + 4x + 5)$.

\begin{enumerate}
    \item Quel est le degré de $Q$. Justifier \hfill \textbf{0.5 + 1 pts}
    \item Donner deux racines évidentes de $Q$. \hfill \textbf{1 pts}
    \item Recopier et compléter le tableau de signe de $-x^2 + 4x + 5$ ci-dessous. \hfill \textbf{1.5 pts}

    \vspace{0.2cm}

    \begin{center}
    \begin{tabular}{|c|cc c c|}
    \hline
    $x$ & $-\infty$ & & & $+\infty$ \\ \hline
    $-x^2 + 4x + 5$ & & $0$ & $0$ & \\ \hline
    \end{tabular}
    \end{center}

    \vspace{0.2cm}

    \item Recopier et compléter le tableau de signe de $Q(x)$ ci-dessous. \hfill \textbf{2 pts}

    \vspace{0.2cm}

    \begin{center}
    \begin{tabular}{|c|cccccc|}
    \hline
    $x$ & $-\infty$ & $-3$ & $-1$ & $1$ & $5$ & $+\infty$ \\ \hline
    $x + 3$ & & & & & & \\ \hline
    $x - 1$ & & & & & & \\ \hline
    $-x^2 + 4x + 5$ & & & & & & \\ \hline
    $Q(x)$ & & & & & & \\ \hline
    \end{tabular}
    \end{center}

    \vspace{0.2cm}

    \item En déduire de ce qui précède la solution de : \\
    a) l'équation $Q(x) = 0$. \quad b) L'inéquation $Q(x) > 0$. \hfill \textbf{1 + 1 pts}
    \item Déterminer la forme réduite du polynôme $Q(x)$. \hfill \textbf{1 pts}
\end{enumerate}

\fbox{\textbf{\exo}(11 points)}

\vspace{0.2cm}

\noindent 1) On considère le polynômes $P$ définie par: $P(x) = -x^3 + 5x^2 + 3x - 15$.
\begin{enumerate}[label=\alph*)]
    \item Calculer $P(\sqrt{3})$ et $P(5)$. \hfill \textbf{1 + 0.5 pts}
    \item Déterminer $a$ et $b$ tel que : $P(x) = (x - \sqrt{3})(x - 5)(ax + b)$. \hfill \textbf{1+1 pts}
    \item Résoudre dans $\mathbb{R}$ les équations : $P(x) = 0$ et $P\left(x + \sqrt{3}\right) = 0$. \hfill \textbf{1 + 1 pts}
    \item Résoudre dans $\mathbb{R}$ les inéquations : $P(x) \geq 0$ et $P(x) < -15$. \hfill \textbf{1 + 1.5 pts}
\end{enumerate}

\noindent 2) On donne les applications de $\mathbb{R}$ dans $\mathbb{R}$ suivantes : $f(x) = \frac{-x}{x^2 + 3}$ et $g(x) = 1 - x^2$.
\begin{enumerate}[label=\alph*)]
    \item Calculer $fof(-1)$ et $gof(2)$. \hfill \textbf{0.5 + 0.5 pts}
    \item Déterminer les expressions $fog(x)$ et de $gog(x)$ \hfill \textbf{1 + 1 pts}
\end{enumerate}

\fbox{\textbf{\exo}(14 points)}
\begin{enumerate}
    \item On donne le polynôme : $P(x) = ax^3 + bx^2 - 5x + c$ où $a$ ; $b$ et $c$ sont des réels non nuls.
    \begin{enumerate}
        \item Quelle est le degré de $P(x)$ ? \hfill (1pt)
        \item Déterminer $a$ ; $b$ et $c$ sachant que $P(-2) = 0$ ; $P(0) = 2$ et $P(1) = 0$. \hfill (2pts)
    \end{enumerate}
    
    \item Dans la suite, on considère que $P(x) = 2x^3 + x^2 - 5x + 2$.
    \begin{enumerate}
        \item Factoriser $P(x)$ en produit de facteurs du premier degré. \hfill (2pts)
        \item Résoudre dans $\mathbb{R}$, l'équation $P(x) = 0$. \hfill (1pt)
        \item Résoudre dans $\mathbb{R}$, l'inéquation $P(x) > 0$. \hfill (2pts)
        \item Déduire de la question 2.b. la solution dans $\mathbb{R}$ de l'équation
        \[ 2(-x + 1)^3 + (-x + 1)^2 - 5(-x + 1) + 2 = 0. \] \hfill (2pts)
    \end{enumerate}
    
    \item On pose : $f(x) = \frac{P(x)}{x^2 - 4}$.
    \begin{enumerate}
        \item Déterminer l'ensemble de définition $D_f$ de $f$. \hfill (1pt)
        \item Simplifier $f(x)$. \hfill (1pt)
        \item Etudier le signe de $f(x)$. \hfill (2pts)
    \end{enumerate}
\end{enumerate}


\fbox{\textbf{\exo}(10 points)}

\noindent $Q(x) = a\,x^3 + b\,x^2 - 25x + 12$

\begin{enumerate}
    \item Trouver $a$ et $b$ tels que $Q(3) = 0$ et $Q(1) = -10$ \hfill \textbf{(2 pts)}
    \item Dans tout le reste de l'exercice, on prend $a = 2$ et $b = 1$.
    \begin{enumerate}
        \item Factoriser complétement $Q(x)$. \hfill \textbf{(2 pts)}
    \end{enumerate}
    \item On pose $P(x) = (-2x + 3)Q(x)$.
\end{enumerate}

\begin{enumerate}
    \item Résoudre dans $\mathbb{R}$ l'équation $P(x) = 0$ et l'inéquation $P(x) \leq 0$ \hfill \textbf{(3 pts)}
    \item On pose $f(x) = \frac{P(x)}{x + 3}$. Etudier le signe de $f(x)$. \hfill \textbf{(3 pts)}
\end{enumerate}

\fbox{\textbf{\exo}(07 points)}
\noindent Soit un polynôme $P$ défini par $P(x) = 2x^3 + 5x^2 - 11x - 14$

\begin{enumerate}
    \item Montrer que 2 est une racine de P puis conclure \hfill (0,75 pts)
    \item Factoriser complètement $P(x)$ \hfill (1,5 pts)
    \item Résoudre $P(x) = 0$ ; $P(x) < 0$ \hfill (2pts)
    \item En déduire les solutions de : $P(2x + 1) = 0$ \hfill (1pt)
    \item On pose $f(x) = \frac{P(x)}{x^2 + x - 6}$
    \begin{enumerate}
    \item Déterminer le domaine de définition $D_f$ de $f$ \hfill (0,5pt)
    \item Simplifier $f(x)$ \hfill (0,5pt)
    \item Résoudre $f(x) \geq 0$ \hfill (0,75pt)
\end{enumerate}
\end{enumerate}

\fbox{\textbf{\exo}(07 points)}

\begin{enumerate}
    \item On donne le polynôme $P(x) = 2x^3 + ax^2 - 18x + b$ ; ou a et b sont réels.\\
    Déterminer a et b sachant que $P(\frac{1}{2}) = 0$ et $P(2) = 0$.\hfill \textbf{(1,5pts)}

    \item Dans la suite, on considère que $P(x) = 2x^3 + 3x^2 - 18x + 8$ .
    \begin{enumerate}
        \item Factoriser complètement $P(x)$\hfill \textbf{(1,5pts)}
        \item Résoudre dans $\mathbb{R}$ ,l'équation: $P(x) = 0$\hfill \textbf{(1 pt)}
        \item Résoudre dans $\mathbb{R}$ ,l'inéquation : $P(x) \leq 0$\hfill \textbf{(1,5 pts)}
    \end{enumerate}

    \item Déduire de la question 2) b) les solutions dans $\mathbb{R}$ de :\\
    (E) : $2(x - 3)^3 + 3(x - 3)^2 - 18(x - 3) + 8 = 0$\hfill \textbf{(1,5pt)}
\end{enumerate}

\fbox{\textbf{\exo}(08 points)}

\noindent Soit le polynôme $\text{P(x)} = 2\,x^3 + ax^2 + bx + 2$ où $a$ et $b$ sont des réels.

\begin{enumerate}
    \item Déterminer les réels a et b sachant que : $\text{P}(-1) = 0$ et $\text{P}(1) = 2$ \hfill \textbf{(01pt)}
    \item On pose $\text{P}(x) = 2x^3 - x^2 - x + 2$
    \begin{enumerate}
        \item Factoriser $\text{P}(x)$. \hfill \textbf{(01pt)}
        \item Résoudre dans $\mathbb{IR}$ l'équation $\text{P}(x) = 0$ \hfill \textbf{(01pt)}
        \item Résoudre dans $\mathbb{IR}$ l'inéquation $\text{P}(x) \geq 0$ \hfill \textbf{(01pt)}
    \end{enumerate}
    \item Soit $\text{f}(x) = \frac{x^3 + 2x^2 - 3}{-2x^2 - 3x + 5}$
    \begin{enumerate}
        \item Déterminer le domaine de définition $\text{D}_f$ de f. \hfill \textbf{(0,75pt)}
        \item Factoriser le numérateur et le dénominateur de f, puis simplifier $\text{f}(x)$. \hfill \textbf{(01,5pt)}
        \item Résoudre dans $\mathbb{IR}$ :
        \begin{itemize}
            \item[$\succ$] $\text{f}(x) = 0$ \hfill \textbf{(0,75pt)}
            \item[$\succ$] $\text{f}(x) \leq 0$ \hfill \textbf{(01pt)}
        \end{itemize}
    \end{enumerate}
\end{enumerate}

\fbox{\textbf{\exo}(10 points)}

\vspace{0.2cm}

\noindent\textbf{\underline{Partie A}}

\vspace{0.2cm}

\noindent On considère le polynôme $P(x) = 2x^3 - 7x^2 + \mathbf{a}x + \mathbf{b}$.

\begin{enumerate}
    \item Déterminer les réels \textbf{a} et \textbf{b} sachant que $P(-1) = 0$ et $P(1) = 6$. \hfill \textbf{1 pt}
    \item Pour la suite on donne $P(x) = 2x^3 - 7x^2 + x + 10$.
    \begin{enumerate}
        \item Vérifier que $P(x)$ est factorisable par $(x + 1)$. \hfill \textbf{1 pt}
        \item Déterminer le polynôme $Q(x)$ tel que $P(x) = (x + 1)Q(x)$. \hfill \textbf{1 pt}
        \item Factoriser le trinôme du second degré $Q(x)$. \hfill \textbf{1 pt}
        \item En déduire une factorisation complète de $P(x)$. \hfill \textbf{0,5 pt}
    \end{enumerate}
    \item Résoudre dans $\mathbb{R}$ l'équation $P(x) = 0$. \hfill \textbf{1 pt}
    \item Résoudre dans $\mathbb{R}$ l'inéquation $P(x) < 0$. \hfill \textbf{1 pt}
\end{enumerate}

\vspace{0.4cm}

\noindent\textbf{\underline{Partie B}}

\vspace{0.2cm}

\noindent On donne la fonction $f$ définie par :
\[ f(x) = \frac{P(x)}{-3x^2 - x + 2} \]

\begin{enumerate}
    \item Montrer que : $-3x^2 - x + 2 = (x + 1)(-3x + 2)$. \hfill \textbf{0,5 pt}
    \item En déduire l'ensemble de définition $D_f$ de la fonction $f$. \hfill \textbf{1 pt}
    \item Simplifier $f(x)$ en utilisant la forme factorisée de $P(x)$. \hfill \textbf{1 pt}
    \item Résoudre dans $D_f$ l'inéquation $f(x) \geq 0$. \hfill \textbf{1 pt}
\end{enumerate}

\end{document}